\documentclass[10pt,a4paper]{amsart}
\usepackage[margin=19mm]{geometry}
\usepackage{amsmath,amssymb,mathtools}
\usepackage[scr=rsfs,cal=euler]{mathalfa}
\usepackage{xfrac}
\usepackage[unicode,hidelinks,bookmarks=false]{hyperref}
\newcommand{\ie}{i.e.\ }
\newcommand{\cf}{cf.\ }
\newcommand{\resp}{respectively\ }
\newcommand{\Subsection}{\subsection}
\providecommand{\nobreakdash}{\nobreak\hbox{-}}
\setlength{\emergencystretch}{2em}
\multlinegap0pt
\usepackage{bm}
\usepackage{bbm}
\usepackage{dsfont}
\usepackage{xspace}
\usepackage{cancel}
\usepackage{mathtools}
\usepackage{amsthm}
\makeatletter
\@ifundefined{definition}{\newtheorem{definition}{Definition}}{}
\@ifundefined{fact}{\newtheorem{fact}[definition]{Fact}}{}
\@ifundefined{proposition}{\newtheorem{proposition}[definition]{Proposition}}{}
\makeatother
\DeclareMathOperator{\Int}{int}
\DeclareMathOperator{\cl}{cl}
\def \M {{\mathcal M}}
\newcommand{\N}{\mathbb{N}}
\newcommand{\bN}{\mathbb N}
\newcommand{\ov}{\overline}
\title[Open cell property in weakly o-minimal structures]{Open cell property in weakly o-minimal structures}
\author{Tomohiro Kawakami, Hiroshi Tanaka}
\date{Source: arXiv:2510.26013v2 (CC BY 4.0)}
\begin{document}
\maketitle

\noindent\emph{Source and licence.} Open cell property in weakly o-minimal structures. Version arXiv:2510.26013v2 (CC BY 4.0), distributed under the Creative Commons Attribution 4.0 International licence, as recorded in the arXiv licence metadata for this exact version and shown on the abstract page https://arxiv.org/abs/2510.26013v2. Full rights notice: licenses/arxiv-2510.26013-CC-BY-4.0.txt.\par\smallskip

\noindent\textit{Editorial scope.} Counted unit: the Proposition whose source label is \texttt{open} in the article (source0.tex lines 526--532, source label \texttt{open}), delivered together with the article's own complete original proof of it (source0.tex lines 535--647, one \texttt{proof} environment of 113 lines). No mathematics is written, repaired or re-derived: statement and proof are the article's own text line for line. Required context retained uncounted: basic (lines 514--519). Every definition, display and label of those lines is retained unchanged. Omitted from this package and not redistributed: 1--513, 520--525, 533--534, 648--1000. Nothing inside a retained excerpt is omitted or altered: every retained body is delivered exactly as the article's lines are written, so no omission receipt of any kind accompanies this package. Citations: the retained excerpts carry no citation command at all, so no bibliography entry of the article is transcribed and no reference is redistributed. Reference adaptation: none was needed and none was made; every reference inside the delivered unit resolves inside it, so no unresolved reference or citation is produced. Printed slips kept literal: no printed word, symbol, hypothesis or number of the counted statement or of its proof was altered. Layout and class: the article's own class is \texttt{amsart} with packages including amsmath, amssymb, amsthm, enumerate, graphicx; the wrapper is the lane's pinned \texttt{amsart} editorial wrapper at 10pt on A4, which restates the theorem-like environments the retained text needs and adds the article's own macro definitions that the retained text needs (\texttt{\textbackslash Int}, \texttt{\textbackslash M}, \texttt{\textbackslash N}, \texttt{\textbackslash bN}, \texttt{\textbackslash cl}, \texttt{\textbackslash ov}), with extra wrapper packages (bm, bbm, dsfont, xspace, cancel, mathtools). The delivered numbering is the unit's own. Assets: no retained span contains a figure environment, a graphics-inclusion command, a PDF-page inclusion, a table or a TikZ picture; no image file is copied into this package, the package declares no asset dependency and no image file is redistributed with it. Licence: version arXiv:2510.26013v2 of this article is distributed under the Creative Commons Attribution 4.0 International licence, as recorded in the arXiv licence metadata for this exact version and shown on the abstract page https://arxiv.org/abs/2510.26013v2. A full rights notice is delivered at licenses/arxiv-2510.26013-CC-BY-4.0.txt. Characters: every retained excerpt is pure ASCII, so the delivered engine, XeLaTeX with its default Latin Modern text font, reports no missing character.
\begin{fact} \label{basic}
	Assume that $\mathcal M = (M, <, \dots)$ is a weakly o-minimal structure with the strong cell decomposition property and $\ov{\mathcal M} = (\ov{M}, <, \dots)$ is its canonical o-minimal extension.
	Let $m, k \in \N_{+}$.
	If $X_1, \dots, X_k \subseteq \ov{M}^m$ are sets definable in $\ov{\M}$,
	then there exists a basic cell decomposition of $\ov{M}^m$ into finitely many basic cells partitioning each of the sets $X_1, \ldots, X_k$.
\end{fact}

\begin{proposition} \label{open}
Suppose that $\mathcal M = (M, <, \ldots)$ is a weakly o-minimal structure with the refined strong cell decomposition property.
Let $n \in \bN_{+}$.
Suppose that  $X$ is a definable open subset of $M^n$.
Then there exists an $\overline{\M}$-definable open set $Y \subseteq \overline{M}^n$ such that
$X = Y \cap M^n$.
\end{proposition}

\begin{proof}
	Take a refined strong cell decomposition of $X$, so that
	\[
	X = C_1 \cup \dots \cup C_k.
	\]
	By Fact~\ref{basic},
	there exists a basic cell decomposition $\mathcal{C}$ of $\overline{M}^n$ partitioning each of the sets $\ov{C_1},\dots, \ov{C_k}$.
	Define
	\[
	\hat{X}
	= \bigcup \left\{ D\in \mathcal{C} : \exists i\ (D \subseteq \overline{C_i}) \right\}
	\;\cup\;
	\bigcup \left\{ D\in \mathcal{C} : \bigl(\forall i\ (D \cap \overline{C_i} = \emptyset)\bigr) \wedge (D \cap M^n = \emptyset) \right\}.
	\]
	Set $Y = \Int_{\overline{M}}(\hat{X})$.
	Since $\mathcal{C}$ is finite, $\hat{X}$ is definable in $\overline{\M}$; hence $Y$ is definable  in $\overline{\M}$ as well.
	We prove that $X = Y \cap M^n$.
	
	By the definition of $\hat{X}$, the inclusion $X \supseteq Y \cap M^n$ is clear.
	
	It remains to show $X \subseteq Y \cap M^n$.
	Assume for a contradiction that there exists $a=(a_1,\dots,a_n)\in X$ such that
	$a \notin \Int_{\overline{M}}(\hat{X})$.
	Since $X$ is open in $M^n$, there exists an open box $U_0 \subseteq X$ in $M^n$ containing $a$.
	Since $a \notin \operatorname{int}_{\overline{M}}(\hat{X})$, there exists a point $b_0 \in \overline{U_0}$ such that $b_0 \notin \hat{X}$.
	Let $U_1$ be an open box in $M^n$ containing $a$ such that
	$\cl_M(U_1) \subseteq U_0$ and $b_0 \notin \overline{U_1}$.
	Then there exists a point $b_1 \in \overline{U_1}$ such that $b_1 \notin \hat{X}$.
	Iterating this argument, there are infinitely many points
	$b_i \in \overline{U_0}$ ($i\in\omega$) with $b_i \notin \hat{X}$.
	Let $B = \{ b_i : i \in \omega \}$.
	Since $\mathcal{C}$ is finite, there exists $C \in \mathcal{C}$ such that infinitely many of the $b_i$ lie in $C$.
	Replacing $B$ by this infinite subset if necessary, we may assume $B \subseteq C$.
	Moreover, since $a \notin \Int_{\overline{M}}(\hat{X})$, for every open box $U \subseteq X$ in $M^n$ containing $a$ there exists $b \in \overline{U}$ with $b \notin \hat{X}$;
	hence we may assume $a \in \cl_{\overline{M}}(C)$.
	Also, $C$ is infinite.
	
	If $C \subseteq \overline{C_i}$ for some $i$, then $B \subseteq \overline{C_i} \subseteq \hat{X}$, a contradiction.
	Thus for every $i$ we have $C \cap \overline{C_i} = \emptyset$, and hence $C \cap X = \emptyset$.
	
	First, consider the case where $C$ is a basic $\langle 0,\dots,0,1 \rangle$-cell.
	Then there exists an open interval $I \subseteq \overline{M}$ such that
	\[
	C = (a_1,\dots,a_{n-1}) \times I.
	\]
	Since $a \in X \subseteq M^n$ and $a \in \cl_{\overline{M}}(C)$, every open box $U$ in $M^n$ containing $a$ satisfies
	$C \cap U \ne \emptyset$, equivalently $(C \cap M^n) \cap U \ne \emptyset$.
	Because $C \cap X = \emptyset$, we have $(C \cap M^n) \cap X = \emptyset$.
	Hence $a \notin \Int_M(X)$, contradicting that $X$ is open in $M^n$.
	
	Next, consider the case where $C$ is a basic $\langle i_1,\dots,i_{n-1},0 \rangle$-cell and $i_1+\dots+i_{n-1} > 0$.
	That is, there exists a basic $\langle i_1,\dots,i_{n-1} \rangle$-cell $\hat{D} \subseteq \overline{M}^{n-1}$, and if we set
	\[
	\{j_1,\ldots,j_k\} = \{\, j \in \{1,\dots,n-1\} : i_j = 1 \,\},
	\qquad j_1 < \dots < j_k,
	\]
	then
	\[
	D = \rho_{j_1,\ldots,j_k}(\hat{D}) \cap M^k
	\]
	is an open strong cell in $M^k$.
	Moreover, there exists a strongly continuous definable function $f$ from $D$ to $M$ or from $D$ to $\overline{M}\setminus M$ such that
	\[
	C = \Gamma\bigl(\overline{f} \circ \rho_{j_1,\ldots,j_k}(\hat{D})\bigr).
	\]
	
	Assume that $f:D\to M$.
	Recall that $C \cap \overline{C_i} = \emptyset$ for every $1 \leq i \leq k$.
	If $C \cap M^n = \emptyset$, $C \subseteq \hat{X}$ by the definition of $\hat{X}$.
	Hence $B \subseteq C \subseteq \hat{X}$, a contradiction.
	
	Assume that $C \cap M^n \ne \emptyset$.
	In the construction of a function-type cell in the basic cells, all cell-defining functions take values in $M$.
	We claim that $a \in \cl_{\overline{M}}(C \cap M^n)$.
	Let $U \subseteq M^n$ be an arbitrary open box containing $a$.
	Since $a \in \cl_{\overline{M}}(C)$, there exists a point $x=(x_1,\dots,x_n) \in \overline{U} \cap C$.
	Set $\tilde{C} = \rho_{j_1,\ldots,j_k}(C)$ and $V = \rho_{j_1,\ldots,j_k}(U)$.
	Then $\tilde{C}$ is open in $\overline{M}^k$, and $V$ is an open box in $M^k$.
	Let $\tilde{x} = (x_{j_1},\dots,x_{j_k})$.
	Then $\tilde{x} \in \overline{V} \cap \tilde{C}$.
	Since $\tilde{C}$ is open in $\overline{M}^k$, there exists an open box $W$ in $M^k$ such that
	$\tilde{x} \in \overline{W} \subseteq \overline{V} \cap \tilde{C}$.
	Choose $\tilde{y} \in W$.
	Since $f:D\to M$, there exists $y \in (C \cap M^n) \cap U$ such that $\rho_{j_1,\ldots,j_k}(y) = \tilde{y}$.
	Therefore $(C \cap M^n) \cap U \ne \emptyset$, and hence $a \in \cl_{\overline{M}}(C \cap M^n)$.
	However, since $C \cap X = \emptyset$, we have $(C \cap M^n) \cap X = \emptyset$, and hence $a \notin \Int_M(X)$.
	This contradicts the assumption that $X$ is open in $M^n$.
	
	Assume instead that $f:D\to \overline{M}\setminus M$.
	Then $C \cap M^n = \emptyset$, and similarly $C \subseteq \hat{X}$.
	Hence $B \subseteq C \subseteq \hat{X}$, a contradiction.
	
	Therefore the function-type case leads to a contradiction.
	
	Finally, consider the case where $C$ is a basic $\langle i_1,\dots,i_{n-1},1 \rangle$-cell and $i_1+\dots+i_{n-1} > 0$.
	As before, there exist a basic $\langle i_1,\dots,i_{n-1} \rangle$-cell $\hat{D} \subseteq \overline{M}^{n-1}$ and an open strong cell $D \subseteq M^k$.
	Moreover, there exist strongly continuous functions $f,g : D \to \overline{M} \cup \{\pm\infty\}$ such that all values of $f$ and $g$ lie in exactly one of
	$\{-\infty\},\ M,\ \overline{M}\setminus M,\ \{\infty\}$, and
	$\overline{f}(x) < \overline{g}(x)$ for any $x \in \overline{D}$, and
	\[
	C = (\overline{f} \circ \rho_{j_1,\ldots,j_k},\ \overline{g} \circ \rho_{j_1,\ldots,j_k})_{\hat{D}}.
	\]
		Recall that $C \cap \overline{C_i} = \emptyset$ for every $1 \leq i \leq k$.
	If $C \cap M^n = \emptyset$, $C \subseteq \hat{X}$ by the definition of $\hat{X}$.
	Hence $B \subseteq C \subseteq \hat{X}$, a contradiction.
	
	
	Assume $C \cap M^n \ne \emptyset$.
	As in the previous case, we have $a \in \cl_{\overline{M}}(C \cap M^n)$.
	Hence every open box $U$ in $M^n$ containing $a$ meets $C \cap M^n$.
	Because $C \cap X = \emptyset$, we have $(C \cap M^n) \cap X = \emptyset$, so $a \notin \Int_M(X)$, contradicting openness of $X$ in $M^n$.
	This completes the proof.
\end{proof}

\end{document}
